% L15a strategy deck: the betting-against-beta (BAB) weighting rule.
% One short deck per rule on the strategy menu, read next to the rule's code file
% (strategies/weighting/bab.jl). Sources: weigh_bab and MyL15aBABRule in bab.jl, MPC.closest_book in
% code/src/mpc/CourseL15a.jl, and the engine settings ls_long = 1.30, ls_short = 0.30, and
% beta_floor = 0.1 in code/src/data/mpc/study.toml. The figure is the right panel of
% figs/Fig-L15a-BAB-Tilt.pdf. Build with make in this folder: LaTeX runs from strategies/, so paths are relative to it.
% November 27, 2026: first build, the pilot of the strategy decks. The figure trim (229bp) keeps the right panel.
\documentclass[aspectratio=169,11pt]{beamer}
\usepackage{vnslides}

\newcommand{\repo}{https://github.com/varnerlab/CHEME-5660-CourseRepository-Fall-2026/blob/main/lectures/week-15/L15a}

\title{\texorpdfstring{{\vnheavy L15a}}{L15a}: Betting Against Beta (BAB)}
\subtitle{Weighting Rule, CHEME 5660 Strategy Menu, Cornell University}
\author{Copyright Jeffrey Varner 2026. All Rights Reserved}

\begin{document}

\vntitlepage

\begin{frame}{Disclaimer and Risks}
  \vspace{0.6em}
  \textbf{\ink{Training only.}} This content is for training and informational
  purposes only. The teaching team does not make, give, or endorse any offer or
  solicitation to buy or sell securities or derivative products, or provide
  investment or trading advice or strategy.

  \vspace{1.1em}
  \textbf{\ink{Risk.}} Trading involves risk and is generally not recommended for
  those with limited resources, experience, or a low risk tolerance. Past
  performance does not guarantee future success.

  \vspace{1.1em}
  \textbf{\ink{Responsibility.}} You are responsible for your investment decisions
  based on your financial circumstances, objectives, risk tolerance, and
  liquidity needs.
\end{frame}

\begin{frame}{The Idea}
  \vspace{0.3em}
  Previously, we met \href{https://www.aqr.com/Our-Firm/About-Us}{AQR Capital Management}, where
  Andrea Frazzini and Lasse H. Pedersen proposed the
  \href{https://www.aqr.com/Insights/Research/Journal-Article/Betting-Against-Beta}{betting-against-beta
  portfolio} in 2014. The BAB weighting rule splits the basket into two sides:
  \begin{itemize}\setlength{\itemsep}{0.5em}
    \item \hilite{Long side}: buy the picked stocks with a low market sensitivity $\beta_{i}$.
    \item \hilite{Short side}: sell short the picked stocks with a high $\beta_{i}$.
    \item \hilite{Sizing}: scale each side to a market sensitivity of one, so that the portfolio's
      $\beta$ is zero before the trading limits.
  \end{itemize}

  \vspace{0.5em}
  The basket comes from the picking rule.

  \vspace{0.5em}
  With fixed estimates, each $\beta_{i}$ keeps its value from the 2016 to 2024 fit of the single index model.
  With updated estimates, the engine updates $\beta_{i}$ at every close. The \href{\repo/strategies/slides/CHEME-5660-L15a-Strategy-Sector-Factor-Fall-2026.pdf}{sector factor estimates}
  change only the expected growth rates, which reach this rule through the basket and the decision days.
\end{frame}

\begin{frame}{The BAB Weights}
  \vspace{0.3em}
  Rank the $K$ stocks of the basket by $\beta_{i}$, from $z_{i}=1$ (lowest) to $z_{i}=K$, with
  the middle rank $\bar{z}=(K+1)/2$:
  {\setlength{\jot}{3pt}
  \begin{align*}
    k &= 2\Big/\sum_{i=1}^{K}\left|z_{i}-\bar{z}\right| && \text{each side adds to one}\\
    w_{i}^{L} &= k\,\max\left(\bar{z}-z_{i},\,0\right),\qquad
    w_{i}^{H} = k\,\max\left(z_{i}-\bar{z},\,0\right) && \text{below or above the middle}\\
    \beta_{L} &= \sum_{i=1}^{K}w_{i}^{L}\beta_{i},\qquad \beta_{H}=\sum_{i=1}^{K}w_{i}^{H}\beta_{i}
      && \text{each side's }\beta\\
    \Longrightarrow\quad &\boxed{w_{i}=\frac{w_{i}^{L}}{\max\left(\beta_{L},\beta_{\text{floor}}\right)}
      -\frac{w_{i}^{H}}{\max\left(\beta_{H},\beta_{\text{floor}}\right)}}
  \end{align*}}%
  where $\beta_{\text{floor}}=0.1$, and a negative $w_{i}$ is a short position. When both sides
  have $\beta$ at least $\beta_{\text{floor}}$, the portfolio's $\beta$ is $1-1=0$.
\end{frame}

\begin{frame}{The Weights of Eight Made-Up Stocks}
  \vspace{0.3em}
  Each weight grows with the stock's distance from the middle rank:
  \begin{center}
    \includegraphics[trim={229bp 0 0 0},clip,width=\textwidth,height=0.59\textheight,keepaspectratio]{../figs/Fig-L15a-BAB-Tilt.pdf}
  \end{center}
  The short side is smaller than the long side. Its stocks have the larger $\beta$, so dividing
  by $\beta_{H}$ shrinks its weights more.
\end{frame}

\begin{frame}{Under the Trading Limits}
  \vspace{0.3em}
  At each decision, the engine finds the portfolio whose weights are closest to the BAB weights,
  in the sum of the squared gaps, within these limits:
  \begin{itemize}\setlength{\itemsep}{0.5em}
    \item \hilite{Long and short sides}: the long weights add to at most 130\% of the wealth, and
      the short weights to at most 30\%.
    \item \hilite{Signs}: a stock with a positive BAB weight can only be held long, and a stock
      with a negative BAB weight can only be held short.
    \item \hilite{Shared limits}: the cash stays at or above zero, and the weight, volatility, and
      trading-cost limits of every strategy apply. With the hedge on, so does the market
      sensitivity limit.
  \end{itemize}

  \vspace{0.5em}
  When a limit stops the engine from reaching the BAB weights, the portfolio's $\beta$ can move
  away from zero.
\end{frame}

\begin{frame}{Reading the Code}
  \vspace{0.3em}
  The rule's code file is
  \href{\repo/strategies/weighting/bab.jl}{\texttt{strategies/weighting/bab.jl}}:
  \begin{itemize}\setlength{\itemsep}{0.5em}
    \item \texttt{weigh\_bab(beta; beta\_floor)} computes the BAB weights $w_{i}$, without the
      trading limits.
    \item \texttt{MyL15aBABRule} holds the settings: \texttt{long\_cap = 1.30},
      \texttt{short\_cap = 0.30}, and \texttt{beta\_floor = 0.1}.
    \item \texttt{MPC.rule\_preferences(...)} marks the picked stocks with a score of one, so that
      the rule can find the basket.
    \item \texttt{MPC.allocate(...)} computes the BAB weights of the basket, then calls
      \texttt{MPC.closest\_book(...)}, which applies the trading limits.
  \end{itemize}

  \vspace{0.5em}
  The \href{\repo/strategies/CHEME-5660-L15a-Strategy-Menu-Fall-2026.ipynb}{strategy menu} lists
  every rule, and the
  \href{\repo/engine/CHEME-5660-L15a-Trading-Engine-Fall-2026.ipynb}{trading engine} describes
  the shared limits.
\end{frame}

\end{document}
